\section{Introduction}
\label{sec:introduction}

Many physical systems are intrinsically stochastic. Thermal fluctuations, environmental noise, unresolved degrees of freedom, and random forcing mean that repeated experiments or simulations performed under nominally identical conditions need not produce the same trajectory. The modelling problem is therefore to characterise this distribution for any arbitrary stochastic system. Direct numerical simulation provides samples from this distribution when the governing dynamics are known, but repeated simulation can become expensive when large ensembles are required across many parameter values.

Data-driven generative models provide an alternative. Neural stochastic differential equations represent stochastic dynamics through learned drift and diffusion functions \cite{kidger2021neuralsde}. In previous work, our group used neural SDEs to model the dynamics of experimental magnetic systems \cite{manneschi2025}. However, adversarial training of such models can be difficult, and generated ensembles may fail to reproduce the full variability of the data. Here we instead consider diffusion models, which learn a distribution through a (simpler) denoising objective rather than adversarial optimisation. In a denoising diffusion probabilistic model (DDPM), training trajectories are progressively corrupted with Gaussian noise and a neural network is trained to predict the added noise \cite{sohldickstein2015,ho2020ddpm}. Generation reverses this process to produce new samples from the learned distribution.

We introduce \texttt{diffusion}, an open-source Python package for conditional diffusion modelling of stochastic dynamical systems. Given simulated or measured trajectories and optional conditioning variables, the package learns
\begin{equation}
    p_{\phi}(\mathbf{x}_{0:T}\mid\mathbf{c})
    \approx
    p_{\mathrm{data}}(\mathbf{x}_{0:T}\mid\mathbf{c}).
\end{equation}
Conditioning $\mathbf{c}$ may include physical parameters, initial conditions, or time-dependent driving signals (or, indeed, nothing at all, in which case the model just samples from the learned distribution). The package supports DDPM and reduced-step DDIM sampling \cite{song2021ddim} (for faster inference), scalar and multichannel trajectories (to learn e.g., spectrograms), and dilated one-dimensional convolutional networks for very long time series or those with long-time correlations. Deterministic DDIM sampling may also be differentiated with respect to the conditioning variables, allowing the trained model to be used for gradient-based inverse problems (e.g., given a set of trajectories, estimate the conditioning variables that best explain the observed data).

We validate on four systems of increasing complexity, from a pure stochastic process to an experimental substrate: (i) the Ornstein–Uhlenbeck process, which provides a baseline with closed-form moments and autocorrelation; (ii) a driven stochastic Duffing oscillator with combined white and coloured noise, which is non-linear, stochastic, and potentially chaotic, and is conditioned on the driving field amplitude; (iii) a magnetic domain-wall oscillator~\cite{ababei2021} in the collective-coordinate model, which is phenomenologically Duffing-like, conditioned on the full driving field; (iv) an experimental interconnected nanomagnetic ring array~\cite{vidamour2022}, conditioned on a rotating in-plane field.

Whilst the package is agnostic to the underlying system, the examples illustrate a range of experimental substrates of interest to our group in the field of physical reservoir computing.